\documentclass[12pt]{article}
\usepackage{graphicx}
\usepackage{amsmath}
\usepackage{url}
\usepackage{booktabs}
\usepackage[margin=1in]{geometry}
\usepackage{natbib}
\usepackage{setspace}
\doublespacing

\title{Cross-Border Cannabis Demand and Dispensary Revenue:\\
National Evidence from U.S.\ State Borders}
\author{Mahyar Rezazadeh \\ Babak Sahragard}
\date{August 2026}

\begin{document}
\maketitle

\begin{singlespace}
\begin{abstract}
\noindent We ask three questions about cannabis retail markets in U.S.\ states located near a state where
recreational sales remain illegal. First, does a border premium exist: is revenue per dispensary higher in
counties near a non-legal-state border than in interior counties of the same state, using all border
counties as treatment and all interior counties as control? Using a county-year panel for Colorado and
Washington, we find it does, both using a 70-mile binary border definition (county-clustered $p=0.033$) and
a continuous distance measure (county-clustered $p=0.023$), while total revenue per capita and a
non-cannabis placebo sector show no comparable pattern. Second, is this evidence of a license-cap policy: do
states that restrict entry show the premium on a different margin than states that do not? Testing this
directly with Montana, the one sample state without a statewide numerical cap, we find the reverse pattern
there, a significant per-capita effect and a precise per-dispensary null, consistent with Colorado and
Washington's entry restrictions forcing captive demand onto existing stores rather than new entry. Third, how
is this cross-border demand captured: does the premium scale with the size of the population on the
non-legal side, not just proximity to it? We show it does (county-clustered $p=0.036$), distinguishing a
county near a small non-legal county from a county near a large non-legal metropolitan area. Together, these
three results establish a border premium, connect it to a specific, testable policy mechanism, and show it
responds continuously to the size of the market it draws from.\\[1em]
\textit{JEL classification}: L51, H71, R12, L81\\
\textit{Keywords}: cannabis; cross-border shopping; market entry; license caps; state policy; border economics
\end{abstract}
\end{singlespace}

\section{Introduction}

Recreational marijuana is legal in a growing number of U.S.\ states but remains prohibited in many others. By
the end of 2025, 24 states had legalized recreational marijuana, though Virginia had not yet established a
legal retail market, leaving 23 states with operating recreational cannabis markets and 27 without one. This
generates sharp policy discontinuities across state borders: residents of a non-legal state may reach a
licensed dispensary with a short drive, while those farther from a legal-state border face higher travel
costs and correspondingly less access to legal recreational cannabis.

Such border effects are a standard feature of settings where neighboring jurisdictions tax, regulate, or
prohibit the same good differently. Cross-border responses to price and policy differences have been
documented for cigarettes \citep{Lovenheim2008}, alcohol \citep{AsplundFribergWilander2007,BeattyLarsenSommervoll2009,Stehr2007},
gasoline \citep{BanfiFilippiniHunt2005}, sales taxes \citep{Agrawal2015,BakerJohnsonKueng2021}, and state
lotteries \citep{KnightSchiff2012}. These responses weaken with distance because travel is costly
\citep{ChandraHeadTappata2014} and strengthen with the size of the population able to reach the more
permissive jurisdiction \citep{KnightSchiff2012,ChandraHeadTappata2014}.

We ask whether the same logic applies to recreational cannabis retail, organized around three specific
questions. \citet{HansenMillerWeber2020} show directly that Washington dispensaries near the Oregon border
lost sales once Oregon's own retail market opened in 2015, establishing that consumers cross state lines for
cannabis. \citet{HaoCowan2020} show a related pattern in marijuana-possession arrests on the non-legal side.
Neither paper asks whether a border premium is general across many borders rather than one policy transition,
whether it depends on state licensing policy, or whether it scales with the size of the population it draws
from; these are the three questions this paper takes up in turn.

\textbf{Result 1} asks whether a border premium exists at all: using every county in Colorado and Washington
within 70 miles of a non-legal border as treatment and every other county in those states as control, we
test whether revenue per active dispensary is higher near the border, using both the binary threshold and a
continuous distance measure. \textbf{Result 2} asks whether this reflects license-cap policy specifically:
if the number of dispensaries a state allows cannot adjust to demand, a demand shock should appear as revenue
per existing dispensary rather than as more dispensaries or more revenue per capita. We test this using
Montana, the one state in our data without a statewide numerical cap, as a direct within-sample contrast.
\textbf{Result 3} asks how this demand is captured: whether the premium in Colorado and Washington scales
with the size of the non-legal-side population within reach, not merely with proximity to it, which would
indicate a captive-demand channel rather than an unrelated correlate of being near a state line.

The remainder of the paper is organized as follows. Section 2 reviews the related literature. Section 3
develops a conceptual framework. Section 4 describes the data and Section 5 the empirical strategy. Sections
6 through 8 report Results 1, 2, and 3 in turn. Section 9 reports robustness checks, Section 10 discusses
limitations and policy implications, and Section 11 concludes. Two secondary analyses that are not part of
the paper's three main results, an entry-timing extension and an exploratory entry-density and fiscal-spillover
calculation, are reported in the appendix rather than the main text.

\section{Literature Review}

A large literature documents that consumers shop across borders when neighboring jurisdictions differ in
taxes, prices, or legal access, and that this behavior weakens with distance and strengthens with the size
of the population able to reach the alternative market. \citet{Lovenheim2008} finds substantial cross-border
cigarette purchasing in the United States; \citet{AsplundFribergWilander2007} show Swedish alcohol sales
respond to Norwegian prices with an elasticity that falls with distance from the border;
\citet{BeattyLarsenSommervoll2009} and \citet{Stehr2007} document similar patterns for alcohol and tobacco
near international and interstate borders, respectively; and \citet{BanfiFilippiniHunt2005} find comparable
fuel-tourism effects for gasoline in Swiss border regions. \citet{ChandraHeadTappata2014} model the
cross-border shopping decision explicitly as a trade-off between price savings and travel cost using
Canada--U.S.\ travel data, while \citet{KnightSchiff2012} show, using state lottery sales, that cross-border
responses are largest where dense populations sit close to a more permissive jurisdiction.
\citet{Agrawal2015} and \citet{BakerJohnsonKueng2021} extend this logic to general sales-tax differentials,
showing that local tax-setting and consumer purchasing both respond to a state border's proximity and to
driving time.

Cannabis retail is a natural extension of this literature because state legalization decisions create sharp,
persistent legal discontinuities rather than only price or tax differentials. \citet{HansenMillerWeber2020}
provide the most direct evidence: using the near-universe of Washington's cannabis transactions and Oregon's
2015 retail-market opening as a natural experiment, they show Washington dispensaries near the Oregon border
lost a substantial share of sales once Oregon consumers no longer needed to cross the border.
\citet{HaoCowan2020} study the non-legal side directly, showing that marijuana-possession arrests are higher
in non-legal-state counties bordering a legal state than in counties farther away. Related work documents
spillovers in crime and traffic safety \citep{WuBoatengLang2020,Moyer2026}, further establishing that
state-level cannabis policy differences generate real cross-border behavioral responses.

A separate literature on cannabis market structure is central to Result 2. \citet{HollenbeckUetake2021} show
that Washington's restrictive licensing regime produced substantial market power among incumbent retailers,
and \citet{MacePatelSeegert2020} show that market structure and competition shape the incidence of cannabis
excise taxes. \citet{Yaskewich2022} documents that local governments frequently restrict or prohibit cannabis
retail even where state law permits it, adding a further layer of supply constraint below the state level.
This literature implies that whether cross-border demand appears as more dispensaries or as more revenue at
existing dispensaries is not a fixed feature of the good, but depends on the elasticity of local supply,
directly analogous to the classical relationship between market size and the number of firms a market
supports \citep{BresnahanReiss1991}.

Finally, because cannabis and alcohol may be substitutes or complements, alcohol is not a fully clean
placebo sector for a cannabis-specific claim. \citet{BaggioChongKwon2020} and \citet{AndersonHansenRees2013}
find evidence of substitution between cannabis and alcohol; \citet{MillerSeo2021} estimate that Washington's
legalization reduced both alcohol and cigarette demand enough to affect state tax revenue; but
\citet{WenHockenberryCummings2015} find increased binge drinking following medical marijuana legalization
and \citet{Veligati2020} find no significant net change in alcohol or cigarette sales. \citet{AndersonRees2023},
in a comprehensive review, conclude the relationship is heterogeneous across policies and populations. We
therefore rely on a placebo sector with no plausible consumption relationship to cannabis, general retail
trade, rather than alcohol, for our main specification test in Result 1.

\section{Conceptual Framework}

Consider county $i$ in a legal-cannabis state, at distance $d_i$ from the nearest county in a non-legal
state. Residents of that non-legal county face three options: forgo recreational cannabis, obtain it
illegally, or travel to a legal dispensary in county $i$. Let $P_i$ denote the non-legal-side population
within a feasible travel radius of county $i$. The cross-border demand county $i$'s dispensaries can draw,
$D_i$, is increasing in $P_i$ and decreasing in $d_i$:
\begin{equation}
D_i = D(P_i, d_i), \qquad \partial D_i/\partial P_i > 0, \qquad \partial D_i / \partial d_i < 0.
\end{equation}
This demand adds to local, in-state demand. Let $R_i$ be total dispensary revenue and $N_i$ the number of
active dispensaries. If entry adjusts freely to demand, additional demand should appear partly or wholly as
a larger $N_i$, following the standard relationship between market size and firm counts
\citep{BresnahanReiss1991}. Where entry is constrained by state or local licensing policy
\citep{HollenbeckUetake2021,Yaskewich2022}, $N_i$ cannot adjust, and the same demand should instead raise
revenue per existing dispensary:
\begin{equation}
R_i/N_i = f(D_i(P_i,d_i),\, L_i,\, N_i), \qquad \partial(R_i/N_i)/\partial P_i > 0 \ \text{when } N_i \text{ is supply-constrained},
\end{equation}
where $L_i$ is local demand. This framework maps directly onto our three results. Result 1 tests whether
$R_i/N_i$ is higher where $d_i$ is smaller, using both a binary and a continuous version of $d_i$. Result 2
tests the supply-constraint condition directly: in a state without a binding license cap, $N_i$ is free to
adjust, so the same cross-border demand should appear in revenue per capita rather than revenue per
dispensary. Result 3 tests whether the effect is stronger where $P_i$ is larger, which distinguishes a
genuine captive-demand channel from an unrelated correlate of border proximity, particularly once combined
with a placebo test in a sector without a comparable legal discontinuity at the same border.

\section{Data}
\label{sec:data}

Cannabis revenue comes from each state's own tax administrative data. For Colorado we use the Department of
Revenue's county-level marijuana sales series on the Colorado Information Marketplace
\citep{COMarijuanaSalesByCounty}, which reports monthly medical and recreational sales by county from 2014
onward; we aggregate to calendar years and keep only county-years with all twelve months reported. For
Washington we use the Liquor and Cannabis Board's fiscal-year sales and excise tax reports by county
\citep{WALCBSalesByCounty}, using the retailer-only sales figures in each report. Retail sales began January
1, 2014 in Colorado \citep{COMedLicensees} and July 8, 2014 in Washington \citep{WSLCB2014FirstLicenses}.
Dispensary counts are built from each state's own current licensee registry rather than a geocoded historical
listing: the Colorado Marijuana Enforcement Division's licensee lookup tool \citep{COMedLicensees} and the
Washington Liquor and Cannabis Board's own current retailer list, cross-referenced by county directly as
reported by the licensee record rather than inferred from a geocoded address
\citep{WALCBLicenseeList,WALCBCannabisRetailers2026}, applied as a current snapshot to each county's panel
years, consistent with how the per-store outcome is constructed throughout. Restricting the resulting
county-year panel to observations with at least one open dispensary leaves $n=787$ county-years: 42 of
Colorado's 64 counties and 38 of Washington's 39 counties, the remainder never having an operating dispensary
over the sample period.\footnote{The 22 excluded Colorado counties reflect two overlapping patterns: local
prohibition of retail marijuana under Amendment 64, which gives municipalities and counties authority to ban
retail sales outright (Douglas County, population 383,906 and Colorado's fifth-largest, is excluded on this
basis alone), and small rural counties, seventeen of the twenty-two have populations under 30,000, that never
attracted a licensed dispensary regardless of local policy.} Population comes from the U.S. Census Bureau's
Population Estimates Program, county
population totals, 2023 vintage \citep{CensusPopEst2023}; we use each county's July 1, 2023 estimate
throughout the panel rather than a year-specific series, consistent with treating population as a
slow-moving control rather than a time-varying outcome.

A county is classified as a border county if its Census Bureau internal point, the standard representative
coordinate the Bureau reports for each county rather than a population-weighted centroid, lies within 70
miles of the internal point of the nearest county in a state without an operating recreational retail
market. This classification is fixed for each county over the full sample period; it is built from the most
recent map of state legalization dates available to the project rather than updated year by year as
individual neighboring states legalize, so within-panel changes in border status do not occur by
construction. The continuous distance measure used in Result 1 is the natural log of this same distance in
miles; no specification we report uses inverse distance or distance bins. $P_i$, the potential non-legal-side
market used in Result 3, is the combined population of every county in a non-legal state within 70 miles of
county $i$, not only the nearest one; it uses the same fixed 2023 population vintage and the same non-legal-
state map as the border classification and is therefore also fixed over time.

In Colorado, 6 of the 42 sample counties meet the border definition, bordering Utah, Wyoming, or Nebraska,
and the remaining 36 are interior. Because dispensaries did not open in every county at once, the number of
Colorado counties observed in the panel grows from 25 in 2015 to 41 from 2023 onward; border counties grow
correspondingly from 3 in 2015 to the full set of 6 from 2020 onward, averaging 5.3 border and 30.4 interior
counties per year. In Washington, 6 of the 38 sample counties are border counties, all near Idaho; this count
grows from 4 in 2015 and 2016 to 5 in 2017 and the full set of 6 from 2018 onward, averaging 5.6 border and
30.4 interior counties per year. Table~\ref{tab:sample_construction} summarizes these patterns by year. At
the observation level, the main Colorado-Washington regression sample of $n=787$ county-years splits into
119 border and 668 interior observations (58 and 334 in Colorado, 61 and 334 in Washington).

Montana serves as the contrast case in Result 2 because it imposes no statewide numerical cap on dispensary
licenses. Its adult-use market was established by House Bill 701, signed May 18, 2021 \citep{MTHB701}, with
retail sales beginning January 1, 2022 \citep{MTRevenueCannabis}; retail sales are permitted only in counties
where a majority of voters approved the underlying 2020 ballot initiative. We use the Department of Revenue's
monthly county-level cannabis sales reports \citep{MTRevenueCountySales} and its licensed-dispensary list
\citep{MTDispensaryLocations}. A county-year enters the sample only if all twelve months are disclosed;
Montana's confidentiality rule suppresses county-months with too few licensed operators to protect individual
businesses, and suppression is more common in smaller counties, so the number of counties with complete data
falls from 35 in 2022 and 2023 to 24 and 23 in 2024 and 2025. This is a reporting artifact of the disclosure
rule, not evidence of market contraction. Of the counties with complete data, 11 meet the 70-mile border
definition in 2022 and 2023 (near Idaho, North Dakota, or Wyoming), falling to 9 in 2024 and 8 in 2025 as
smaller border counties are suppressed, averaging 9.8 border and 19.5 interior counties per year over the
four-year sample; at the observation level, the Montana regression sample of $n=117$ county-years splits
into 39 border and 78 interior observations.

The general-retail placebo used in Result 1 uses the Colorado Department of Revenue's county-by-industry
retail sales filings \citep{CDORRetailByCounty}. This source discloses industry-level detail only for the 11
Colorado counties with the largest filing volumes, a count we confirm directly from the underlying filings
rather than from the source's own description; the resulting sample of $n=110$ county-years splits into 30
border and 80 interior observations. The entry-timing extension reported in the appendix draws on a separate,
broader nine-state sample built from each state's own geocoded dispensary opening dates and
retail-legalization date; because that sample is not part of the paper's three main results, we describe it
separately in the appendix where it is used.

\begin{table}[htbp]
\centering
\caption{Sample Construction: Treated and Control Counties by Year}
\label{tab:sample_construction}
\small
\begin{tabular}{lcccl}
\toprule
State & Years & Border counties & Interior counties & Non-legal neighboring states \\
\midrule
Colorado & 2015 & 3 & 22 & UT, WY, NE \\
Colorado & 2016 & 4 & 25 & UT, WY, NE \\
Colorado & 2017--2019 & 5 & 27--30 & UT, WY, NE \\
Colorado & 2020--2025 & 6 & 31--35 & UT, WY, NE \\
Washington & 2015--2016 & 4 & 26--29 & ID \\
Washington & 2017 & 5 & 29 & ID \\
Washington & 2018--2025 & 6 & 31--32 & ID \\
Montana & 2022--2023 & 11 & 24 & ID, ND, WY \\
Montana & 2024 & 9 & 15 & ID, ND, WY \\
Montana & 2025 & 8 & 15 & ID, ND, WY \\
\bottomrule
\end{tabular}
\begin{minipage}{0.95\textwidth}
\vspace{0.3em}
\footnotesize
\textit{Notes:} Border and interior counts are the counties actually observed in each state's revenue panel
in that year, not the full statewide county count. Colorado's and Washington's border definitions are fixed
per county for the full sample; year-to-year changes reflect counties entering the panel as they obtain
their first dispensary, not changes in border status. Montana's decline in 2024--2025 reflects the state's
confidentiality-suppression rule rather than market contraction; see text.
\end{minipage}
\end{table}

\section{Empirical Strategy}
\label{sec:empirical_strategy}

To estimate the effect of proximity to a non-legal-state border on revenue per dispensary in Colorado and
Washington, we use the following baseline specification:
\begin{equation}
\ln(R_{ist}/N_{ist}) = \alpha + \beta_1\,\text{Border}_i + \beta_2\ln(P_i) + \gamma\ln(\text{Population}_{ist}) + \theta_t + \delta_s + \varepsilon_{ist}.
\end{equation}
$R_{ist}$ is total dispensary revenue in county $i$, state $s$, year $t$, and $N_{ist}$ is the number of
dispensaries open in county $i$ as of that year's reporting cutoff, so the outcome is log revenue per active
dispensary, our preferred margin given that the number of dispensaries a state permits is frequently
policy-constrained rather than demand-determined (Section~\ref{sec:result2}). $\text{Border}_i$ is the
treatment variable: the binary 70-mile indicator described in Section~\ref{sec:data}, fixed for each county
over the full sample, so all identifying variation is cross-sectional, between border and interior counties,
rather than within a county over time. $\ln(P_i)$, the log combined population of non-legal-state counties
within 70 miles of county $i$, enters as a control so that $\beta_1$ is not simply picking up border counties
happening to face a larger captive population than interior counties. $\ln(\text{Population}_{ist})$ is each
county's own log population, netting out ordinary market-size effects on revenue per store. $\theta_t$ and
$\delta_s$ are year and state fixed effects, absorbing national demand trends common to both states in a
given year and any fixed level difference in revenue per dispensary between Colorado and Washington.

The coefficient of interest is $\beta_1$: the percentage difference in revenue per dispensary between border
and interior counties within the same state and year, holding fixed each county's own population and its
non-legal-side population exposure. A positive, significant $\beta_1$ indicates that existing dispensaries in
border counties absorb more revenue than an equally populous interior dispensary would, our test of whether a
captive cross-border demand channel operates on the per-store rather than the per-capita margin.

We estimate three variations of this specification. For the \textbf{main sample} (Result~1), we estimate
equation~(3) by OLS on the pooled Colorado-Washington county-year panel, reporting both HC1 and
county-clustered standard errors and treating the clustered estimate as preferred given that each county is
observed repeatedly. For the \textbf{continuous measure}, also Result~1, we replace the binary
$\text{Border}_i$ with $\ln(d_i)$, the log distance in miles to the nearest non-legal-state county, testing
whether the relationship is a smooth distance-decay gradient rather than an artifact of the specific 70-mile
cutoff; a negative coefficient on $\ln(d_i)$ corresponds to the same prediction as a positive coefficient on
the binary indicator. For \textbf{Result~3}, we instead replace $\text{Border}_i$ with $\ln(P_i)$ itself as
the regressor of interest, dropping it as a control, asking whether revenue per store scales with the
\emph{size} of the captive market within reach rather than merely with proximity to it.

For \textbf{robustness}, we re-estimate the main specification (i) excluding Sedgwick County, Colorado, the
single largest outlier in the border group; (ii) separately by state, since Colorado and Washington restrict
entry through different institutional channels (Section~\ref{sec:result2}); (iii) on Montana alone, dropping
$\delta_s$ since Montana is a single state, as a specification test of the license-cap prediction in
Result~2; and (iv) with total revenue per capita and a non-cannabis general-retail sector in Colorado as the
outcome in place of revenue per dispensary, to test whether the border pattern is specific to cannabis retail
rather than a general feature of border-county commerce.

Three threats to this design are worth stating directly. First, because $\text{Border}_i$ never changes
within a county, $\beta_1$ cannot on its own separate a genuine demand effect from a stable, unobserved
difference between border and interior counties. We address this with the continuous-distance specification
and the general-retail placebo, neither of which a stable confound would obviously replicate, and more
directly with the New Mexico natural experiment (Section~\ref{sec:co_nm}), which exploits a within-county
change in border status around an external policy date rather than cross-sectional comparison alone. Second,
reverse causality is unlikely given that dispensary siting and state borders are both fixed well before the
revenue outcomes we observe, though we cannot rule out that unobserved local retail-market conditions jointly
shape both $N_i$ and a county's border status. Third, our core sample is only two states with usable
administrative revenue data, so inference rests on very few state clusters; we report county-clustered rather
than state-clustered standard errors for this reason, and treat the Montana and New Mexico results as
independent, out-of-sample checks of the same mechanism rather than as adding statistical power to the main
estimate.

\section{Result 1: Does a Border Premium Exist?}
\label{sec:result1}

\begin{table}[htbp]
\centering
\caption{Result 1: The Border Premium in Revenue per Dispensary}
\label{tab:result1}
\small
\begin{tabular}{p{6.6cm}cccc}
\toprule
Specification & Coef. & SE (HC1) & SE (cl.) & $p$ (cl.) \\
\midrule
\multicolumn{5}{l}{\textit{Binary: border ($\le$70mi) vs.\ interior}} \\
No controls & $+0.382$ & $0.099$ & $0.219$ & $0.080$ \\
Full controls$^{a}$ & $+0.341$ & $0.075$ & $0.161$ & $0.033$ \\
\quad excl.\ Sedgwick County, CO & $+0.302$ & $0.068$ & $0.147$ & $0.040$ \\
Colorado only & $+0.552$ & -- & $0.199$ & $0.005$ \\
Washington only & $+0.233$ & -- & $0.243$ & $0.334$ \\
\midrule
\multicolumn{5}{l}{\textit{Continuous: log(distance to nonlegal border, miles)}} \\
Full controls$^{a}$ & $-0.215$ & $0.048$ & $0.095$ & $0.023$ \\
\quad excl.\ Sedgwick County, CO & $-0.168$ & $0.048$ & $0.092$ & $0.067$ \\
\bottomrule
\end{tabular}
\begin{minipage}{0.95\textwidth}
\vspace{0.3em}
\footnotesize
\textit{Notes:} Outcome is log(revenue per active dispensary). County-year panel, Colorado and Washington,
$n=787$ (80 counties: 42 Colorado, 38 Washington; 12 border counties: 6 Colorado, 6 Washington; 68 interior).
$^{a}$Full controls: log(population), year fixed effects, state fixed effects. Both the binary and
continuous measures are significant under county-clustered inference; a county closer to a nonlegal border
earns significantly more revenue per dispensary, and the relationship holds using a continuous distance
measure rather than only the 70-mile cutoff.
\end{minipage}
\end{table}


Counties within 70 miles of a non-legal border earn significantly more revenue per dispensary than interior
counties: 41 percent more with the full control set (county-clustered $p=0.033$), robust to excluding the
largest outlier county (Sedgwick County, Colorado). The result holds within Colorado alone
($p=0.005$), though the Washington-specific estimate is directionally consistent but not itself significant
under clustering ($p=0.334$) once Washington's dispensary counts are measured correctly rather than from an
incomplete registry. The continuous distance measure tells the same story and is, if anything, the more
robust of the two: it is significant under both HC1 and county-clustered inference ($p=0.023$), unlike a
simple threshold specification, which is exactly what a genuine distance-decay relationship should look like
rather than an artifact of the 70-mile cutoff.

\begin{table}[htbp]
\centering
\caption{Result 1 Robustness: Which Margin Responds, and a Non-Cannabis Placebo}
\label{tab:result1_placebo}
\small
\begin{tabular}{p{7.4cm}cc}
\toprule
Outcome & Coefficient (border effect) & $p$-value (cl.) \\
\midrule
log(revenue per active dispensary) & $+0.341$ & $0.033$ \\
log(revenue per capita) & $+0.046$ & $0.833$ \\
\midrule
\multicolumn{3}{l}{\textit{Placebo: general retail sales per capita, Colorado, 2016--2025}} \\
Border, full controls & $-0.006$ & $0.947$ \\
\bottomrule
\end{tabular}
\begin{minipage}{0.95\textwidth}
\vspace{0.3em}
\footnotesize
\textit{Notes:} The revenue-per-dispensary premium in Result 1 does not appear in total per-capita revenue,
consistent with the license-cap mechanism in Result 2. The general-retail placebo uses Colorado Department
of Revenue county-industry filings for the 11 counties with disclosed industry detail
\citep{CDORRetailByCounty}, $n=110$ (30 border, 80 interior); the precise null supports a cannabis-specific
rather than a generic border-commerce interpretation of Result 1.
\end{minipage}
\end{table}


The premium is specific to revenue per dispensary. Total revenue per capita shows no comparable border
pattern, and a general-retail placebo in Colorado, a sector without a legal discontinuity at the same
border, returns a precise null ($p=0.947$), arguing against a general border-commerce explanation for
Result 1.

\subsection{A Direct Test: Colorado's Border with New Mexico}
\label{sec:co_nm}

The comparisons above are cross-sectional: border status never changes for a given county, so they cannot
rule out a stable, unobserved difference between border and interior counties that happens to look like a
demand effect. New Mexico's own retail market, however, opened on April 1, 2022, partway through our panel,
switching six Colorado counties (Archuleta, Conejos, Costilla, La Plata, Las Animas, and Montezuma) from
bordering a non-legal state to bordering a legal one. If the premium in Table~\ref{tab:result1} reflects
genuine captive cross-border demand, it should fall in these counties once New Mexico's own market can
absorb that demand, exactly the mechanism \citet{HansenMillerWeber2020} document for Washington and Oregon,
but on a border and event they do not study.

\begin{table}[htbp]
\centering
\caption{Natural Experiment: Colorado's Border with New Mexico}
\label{tab:co_nm}
\small
\begin{tabular}{lcccc}
\toprule
Specification & Coef.\ (NM-border $\times$ post-2022) & SE (HC1) & SE (cl.) & $p$ (cl.) \\
\midrule
Difference-in-differences$^{a}$ & $-0.487$ & $0.106$ & $0.123$ & $<0.001$ \\
\bottomrule
\end{tabular}
\begin{minipage}{0.95\textwidth}
\vspace{0.3em}
\footnotesize
\textit{Notes:} Outcome is log(revenue per active dispensary). Colorado county-year panel, 2014--2025, $n=352$
(2022, the transition year in which New Mexico's market opened partway through, is dropped). Treatment is the
6 Colorado counties within 70 miles of New Mexico (Archuleta, Conejos, Costilla, La Plata, Las Animas,
Montezuma); controls include the 5 Colorado counties bordering a state that remained non-legal throughout the
sample (Utah, Wyoming, Oklahoma, Nebraska, Kansas) and the 31 Colorado interior counties, plus
log(population) and year fixed effects. $^{a}$Post $=1$ for 2023 onward, after New Mexico's April 1, 2022
retail opening. The implied change in revenue per dispensary in New Mexico-border counties after New Mexico's
own market opens is $-38.6$ percent.
\end{minipage}
\end{table}


\begin{figure}[htbp]
\centering
\includegraphics[width=0.75\textwidth]{figures/co_nm_eventstudy.png}
\caption{Event-study path for Colorado counties bordering New Mexico, coefficient relative to 2021, around
New Mexico's April 1, 2022 retail opening (event year 0). Shaded regions are 95\% confidence intervals using
county-clustered standard errors.}
\label{fig:co_nm}
\end{figure}

Table~\ref{tab:co_nm} reports the result. Revenue per dispensary in New Mexico-border counties falls 38.6
percent relative to Colorado's other border and interior counties after New Mexico opens its own market
(2023 onward, with 2022 dropped as a transition year; county-clustered $p<0.001$). Figure~\ref{fig:co_nm}
shows the same pattern as an event-study path against two separate comparison groups: relative to Colorado's
other non-legal-state-border counties, which holds ``being a border county'' fixed and isolates only New
Mexico's change in legal status, the gap is not significantly different from zero in any pre-period year, and
it opens exactly at the 2022 switch and widens monotonically through 2025. This is the pattern a genuine
demand-switching mechanism predicts and a stable, unobserved difference between border and interior counties
could not produce on its own.

We attempted the same design in reverse, replicating \citet{HansenMillerWeber2020}'s own setting, Washington
counties bordering Oregon around Oregon's October 2015 retail opening, using our county-level annual data.
That replication does not cleanly confirm their result; we report it in Appendix~\ref{app:wa_or} rather than
omit it, and view the Colorado-New Mexico result above, not the Washington-Oregon attempt, as this paper's
genuine causal evidence for the captive-demand mechanism.

\section{Result 2: Evidence of License-Cap Policy}
\label{sec:result2}

Colorado and Washington both constrain dispensary entry, though through different institutional channels.
Washington capped the statewide number of retail licenses at 334 from the outset, allocating them by lottery
in oversubscribed jurisdictions \citep{WSLCB2014LotteryPR}. Colorado imposes no statewide cap, but Amendment
64 gives local governments full authority to prohibit retail marijuana entirely or to regulate its number,
and most Colorado municipalities have exercised this option to ban it outright. Under either channel, $N_i$
in our conceptual framework cannot freely expand, so the framework predicts the demand response identified
in Result 1 should appear in revenue per store rather than revenue per capita, which Table~\ref{tab:result2}
already shows for Colorado and Washington.

\begin{table}[htbp]
\centering
\caption{Result 2: Evidence of License-Cap Policy, Capped vs.\ Uncapped States}
\label{tab:result2}
\small
\begin{tabular}{lcc}
\toprule
Outcome & Coefficient (border effect) & $p$-value (county-cl.) \\
\midrule
\multicolumn{3}{l}{\textit{Colorado + Washington (license caps bind)}} \\
log(revenue per dispensary) & $+0.341$ & $0.033$ \\
log(revenue per capita) & $+0.046$ & $0.833$ \\
\midrule
\multicolumn{3}{l}{\textit{Montana (no statewide numerical cap)}} \\
log(revenue per capita), 2022--2025 panel & $+0.654$ & $0.036$ \\
log(revenue per dispensary), 2025 cross-section & $+0.012$ & $0.958^{a}$ \\
\bottomrule
\end{tabular}
\begin{minipage}{0.95\textwidth}
\vspace{0.3em}
\footnotesize
\textit{Notes:} $^{a}$HC1 standard error; too few Montana border counties ($n=23$) for county-clustered
inference. Colorado and Washington constrain entry (Washington through a statewide license cap allocated by
lottery; Colorado through local governments' authority under Amendment 64 to ban or limit retail marijuana).
Montana imposes no equivalent statewide cap. Where entry is capped, the border effect appears in revenue per
dispensary and not per capita; where entry is not capped, it appears in revenue per capita and not per
dispensary, consistent with a license-cap explanation for why the border premium in Result 1 shows up on the
intensive rather than the extensive margin.
\end{minipage}
\end{table}


Montana provides a direct test of whether this is really about license-cap policy, since it is the one
sample state without a statewide numerical cap. If the license-cap channel is correct, Montana's border
effect should appear in revenue per capita rather than revenue per dispensary, the reverse of Colorado and
Washington. It does: a significant per-capita effect ($p=0.036$) alongside a precise per-dispensary null
($p=0.958$). This reversal, occurring in exactly the state whose licensing regime differs in exactly the
relevant way, is considerably stronger evidence for a license-cap mechanism than the Colorado-Washington
result alone could provide, since it is a genuine out-of-sample prediction rather than a pattern fitted after
the fact.

\section{Result 3: How Is Cross-Border Demand Captured?}
\label{sec:result3}

Results 1 and 2 establish that a border premium exists and that it responds to license-cap policy in the
predicted way, but neither speaks to whether the premium is genuinely about the \emph{size} of the captive
market on the non-legal side, as opposed to some other correlate of being close to a state line. Result 3
tests this directly using $P_i$, the combined population of all non-legal-state counties within 70 miles of
county $i$, not merely the nearest one.

\begin{table}[htbp]
\centering
\caption{Result 3: Does the Premium Scale with the Nonlegal-Side Population?}
\label{tab:result3}
\small
\begin{tabular}{p{7.4cm}cccc}
\toprule
Specification & Coef. & SE (HC1) & SE (cl.) & $p$ (cl.) \\
\midrule
log($P_i$: captive nonlegal population, 70mi) + full controls & $+0.030$ & $0.007$ & $0.014$ & $0.036$ \\
\quad excl.\ Sedgwick County, CO & $+0.024$ & $0.007$ & $0.014$ & $0.078$ \\
\bottomrule
\end{tabular}
\begin{minipage}{0.95\textwidth}
\vspace{0.3em}
\footnotesize
\textit{Notes:} Outcome is log(revenue per active dispensary), $n=787$. $P_i$ sums the population of
\emph{all} nonlegal-state counties within 70 miles of county $i$, not only the nearest one, so it
distinguishes a county next to a small nonlegal county from a county next to a large nonlegal metropolitan
area, something the distance measures in Result 1 cannot do. Only 10 of the 80 Colorado and Washington
counties in this panel have any captive nonlegal population within 70 miles at all, so this result reweights
the same identifying set of border counties from Result 1 by market size rather than adding new geographic
variation; combined with Result 1's continuous-distance finding, it indicates the premium responds
continuously to both how close a county is to a nonlegal border and how large the population on the other
side is, not only to a binary border classification.
\end{minipage}
\end{table}


Revenue per dispensary rises with $P_i$, and this result survives county-clustered inference ($p=0.036$).
We view this, together with the continuous-distance result in Result 1, as the strongest evidence that
Result 1's border premium reflects a genuine captive cross-border demand channel: a county next to one small
non-legal county and a county next to a large non-legal metropolitan area are treated identically by a
simple border threshold, but only the latter shows the larger revenue response Result 3 predicts.

\section{Robustness}

Result 1 is robust to excluding the single largest outlier county and to state-by-state estimation. The
general-retail placebo in Section~\ref{sec:result1} addresses the possibility that border counties simply
transact more retail business for reasons unrelated to cannabis policy. Result 3's captive-population measure
addresses the concern that the 70-mile threshold in Result 1 is an arbitrary cutoff, by showing a graded
response to market size rather than only a discontinuous jump at 70 miles. We note directly that the
Washington-specific estimate in Result 1 is not significant under clustering once Washington's dispensary
counts are corrected, a genuine weakening relative to an earlier version of this analysis that relied on an
incomplete dispensary registry; we report the corrected, weaker estimate rather than the more favorable
uncorrected one. The New Mexico natural experiment in Section~\ref{sec:co_nm} addresses a different concern
than these robustness checks: rather than showing the cross-sectional premium survives alternative
specifications, it shows the premium responds to an actual change in a neighboring state's legal status,
which a stable, unobserved difference between border and interior counties cannot explain.

\section{Discussion}

These results imply a measurable fiscal externality between neighboring states with different cannabis
policies: some revenue and associated retail activity that a non-legal state might otherwise capture by
legalizing and taxing recreational cannabis itself instead accrues to legal-state border counties. They also
imply that licensing policy in border counties is not a purely local decision. If border dispensaries serve
a substantially larger effective market than local population alone would suggest, licensing caps set with
only local population in mind may constrain supply well below the level that actual demand, inclusive of
cross-border shoppers, would otherwise support, a concern connected directly to the market-power findings of
\citet{HollenbeckUetake2021}. We do not observe individual consumers' state of residence, so most of our
evidence on the demand channel is indirect, built from distance, population-exposure, and placebo patterns
rather than transaction-level origin data; the New Mexico natural experiment in Section~\ref{sec:co_nm} is
the exception, since it identifies the demand channel from an actual change in a neighboring state's legal
status rather than from cross-sectional comparison alone. Our core revenue analysis for Results 1 and 3 still
rests on only two states with usable administrative data, with Montana serving as a single, informative but
not conclusive, institutional contrast for Result 2.

\section{Conclusion}

We ask three questions about cannabis retail markets near U.S.\ state borders. First, does a border premium
exist: yes, revenue per dispensary is significantly higher in border counties than interior counties, using
both a binary and a continuous distance measure, and the premium does not appear in a non-cannabis placebo
sector; a natural experiment around New Mexico's 2022 legalization shows the premium falls 38.6 percent in
Colorado's New Mexico-border counties once New Mexico's own market can absorb the demand, direct causal
evidence for the mechanism rather than only a cross-sectional association. Second, is this evidence of
license-cap policy: yes, the premium appears on the per-dispensary margin
in Colorado and Washington, which restrict entry, and reverses to the per-capita margin in Montana, which
does not, exactly as a license-cap mechanism predicts. Third, how is this demand captured: the premium scales
with the size of the non-legal-side population within reach, not merely with proximity to it, indicating a
genuine captive-demand channel. Together, these three results extend the cross-border cannabis literature
\citep{HansenMillerWeber2020,HaoCowan2020} from a single border and a single policy transition to a general,
multi-state, and mechanism-tested pattern, and connect it directly to a specific, testable feature of state
cannabis policy.

\appendix

\section{Entry-Timing Extension}
\label{app:entry_timing}

As a secondary analysis outside the paper's three main results, we ask whether entry \emph{speed}, not just
revenue, responds to border proximity, using each dispensary's own opening date across a broader nine-state
sample (California, Colorado, Illinois, Massachusetts, Maine, New York, Ohio, Oregon, and Washington; 479
counties) and each state's own retail-legalization date. We estimate Cox proportional-hazards models
stratified by state for the time from a county's own state retail launch to its first dispensary, and from
its first to its second.

\begin{table}[htbp]
\centering
\caption{Extension: Entry Speed Near the Border (corrected)}
\label{tab:ae_hazard}
\small
\begin{tabular}{lccc}
\toprule
Entrant rank & Hazard ratio, log(distance) & $p$-value & Events / at-risk \\
\midrule
1st (state launch $\to$ 1st dispensary) & $0.93$ & $0.540$ & 271 / 479 \\
2nd (1st $\to$ 2nd dispensary) & $0.61$ & $0.001$ & 211 / 271 \\
3rd (2nd $\to$ 3rd dispensary) & $0.74$ & $0.090$ & 174 / 211 \\
4th (3rd $\to$ 4th dispensary) & $0.84$ & $0.356$ & 142 / 174 \\
\bottomrule
\end{tabular}
\begin{minipage}{0.95\textwidth}
\vspace{0.3em}
\footnotesize
\textit{Notes:} Cox proportional-hazards models stratified by state, controlling log(population); nine-state
entry-density universe, distinct from the two-state revenue panel used elsewhere in this paper.
\textbf{Corrected from the original version of this table}: every Colorado and Washington record in the
underlying geocoded dispensary registry shares an identical placeholder open date (2015-01-01, tagged
\texttt{date\_confidence=baseline\_assumed} in the source file), not an observed opening date. The original
table included these records, producing degenerate zero-duration entry gaps for 58 percent of CO/WA counties
(versus 18 percent pooling all nine states) and inflating the at-risk counts at every rank. This corrected
version excludes all placeholder-dated records, leaving seven states with genuinely observed entry dates
(California, Illinois, Massachusetts, Maine, New York, Ohio, Oregon; 479 counties, 271 with a first entrant).
The second-entrant result is not an artifact of the excluded records: it survives, and its precision improves,
on the clean sample (hazard ratio moves from 0.70, $p=0.003$, on the contaminated 346-county sample to 0.61,
$p=0.001$, on this 271-county corrected sample). A hazard ratio below 1 means faster entry closer to a
non-legal border. The effect remains significant only for the second entrant, weakening at the third
(marginal) and vanishing by the fourth, consistent with a genuine second-entrant phenomenon rather than a
generic proximity-always-speeds-entry pattern; we report it as suggestive secondary evidence, not a central
result.
\end{minipage}
\end{table}


There is no relationship between border proximity and how quickly a county gets its first dispensary
(hazard ratio $0.93$, $p=0.54$), consistent with regulatory and capital-raising delay, not demand, governing
first entry. A county's second dispensary, however, follows significantly faster the closer that county is
to a non-legal border (hazard ratio $0.61$, $p=0.001$), weakening to marginal significance by the third
entrant ($p=0.09$) and vanishing by the fourth ($p=0.36$). A classical cross-sectional entry-threshold
comparison does not show a corresponding shift in the population needed to support a second store. This
nine-state sample was built from a geocoded dispensary registry that, for Washington specifically, we now
know understates the true number of licensed dispensaries; more consequentially for this appendix, every
Colorado and Washington record in that registry carries an identical placeholder open date (2015-01-01) rather
than an observed one, so those two states' counties remain in the population at risk throughout but
contribute no observed entry event to the corrected estimates above -- the second-entrant result is estimated
entirely from the other seven states' genuinely dated records (California, Illinois, Massachusetts, Maine,
New York, Ohio, Oregon), and we no longer report which subset of those seven drives significance without
re-verifying that claim on the corrected sample. We have not re-audited every state's registry in this sample
for the Washington-style undercount issue, so the entry-timing results in this appendix should be read with
that additional caveat in mind, beyond the identification limitations already noted here.

\section{Robustness Check: Replicating Hansen, Miller, and Weber (2020) with County-Level Annual Data}
\label{app:wa_or}

As a robustness check on the New Mexico natural experiment in Section~\ref{sec:co_nm}, we attempted the
mirror-image test in \citet{HansenMillerWeber2020}'s own setting: Washington counties bordering Oregon around
Oregon's October 1, 2015 retail opening, using our own county-level annual revenue and dispensary
construction rather than their monthly transaction data. We compare Washington's 13 Oregon-border counties
(within 70 miles of Oregon) to Washington's 4 Idaho-border counties and its remaining interior counties,
using fiscal years 2015 through 2025 (FY2015 is the only fully pre-legalization year available).

\begin{table}[htbp]
\centering
\caption{Replication Attempt: Washington's Border with Oregon}
\label{tab:wa_or}
\small
\begin{tabular}{lcccc}
\toprule
Specification & Coef. & SE (HC1) & SE (cl.) & $p$ (cl.) \\
\midrule
Transition year only (FY2016 vs.\ FY2015) & $-0.545$ & $0.329$ & $0.336$ & $0.105$ \\
DiD, post $=$ FY2016 onward & $+0.136$ & $0.381$ & $0.432$ & $0.753$ \\
DiD, FY2016 dropped, post $=$ FY2017 onward & $+0.161$ & $0.380$ & $0.443$ & $0.716$ \\
\bottomrule
\end{tabular}
\begin{minipage}{0.95\textwidth}
\vspace{0.3em}
\footnotesize
\textit{Notes:} Outcome is log(revenue per active dispensary). Washington county-year panel, fiscal years
2015--2025, $n=395$ (2015 is the only fully pre-legalization year). Treatment is the 13 Washington counties
within 70 miles of Oregon; controls include the 4 Idaho-border counties, the remaining interior counties,
log(population), and fiscal-year fixed effects. The transition-year estimate is directionally consistent with
\citet{HansenMillerWeber2020} but only marginally significant; the longer-run DiD estimates are positive and
insignificant, the opposite of sustained demand suppression. See Appendix~\ref{app:wa_or} for discussion.
\end{minipage}
\end{table}


\begin{figure}[htbp]
\centering
\includegraphics[width=0.75\textwidth]{figures/wa_or_eventstudy_corrected.png}
\caption{Event-study path for Washington counties bordering Oregon, coefficient relative to fiscal year 2015,
around Oregon's October 1, 2015 retail opening (inside fiscal year 2016). Shaded region is the 95\%
confidence interval using county-clustered standard errors.}
\label{fig:wa_or}
\end{figure}

This replication does not cleanly confirm \citet{HansenMillerWeber2020}'s result. The immediate transition
year (fiscal year 2016, which is roughly three pre-legalization and nine post-legalization months) shows a
directionally consistent dip of similar sign to their finding, though only marginally significant under
clustering ($p=0.10$), but the coefficient reverses to significantly positive from fiscal year 2017 onward and
remains elevated through 2025 (Figure~\ref{fig:wa_or}), the opposite of the sustained suppression a genuine
demand-switching effect would predict. We attribute this to two features of our design rather than to the
underlying mechanism being wrong: our Oregon-border county group is dominated by large, fast-growing
metropolitan counties (Clark County, population 521,150, directly across the Columbia River from Portland;
Thurston, population 299,003; and Benton, population 215,219) whose revenue trends plausibly reflect general
population and urban growth rather than the 2015 switch specifically, and only one full pre-period fiscal
year is available, which is not enough to check for pre-existing divergence between the Oregon-border and
comparison counties before treatment. \citet{HansenMillerWeber2020} avoid both problems with store-level
monthly transaction data and a difference-in-discontinuities design; our coarser annual, county-level data are
simply not suited to this particular replication. We report this null result rather than omit it, and rely on
the New Mexico natural experiment in Section~\ref{sec:co_nm}, whose treated counties are small and rural
rather than large and rapidly urbanizing, for this paper's causal evidence on the demand-switching mechanism.

\section{Exploratory Analysis: Entry Density and Fiscal Spillover}
\label{app:welfare}

As a further secondary analysis, we tested whether the license-cap mechanism in Result 2 also appears on the
entry-density margin directly, that is, whether Colorado and Washington border counties have systematically
fewer dispensaries per capita than interior counties, which would be the textbook signature of a binding
entry restriction. Using the corrected Washington dispensary data described in Section~\ref{sec:data}, this
result does not hold: the Colorado-Washington border coefficient on dispensaries per capita is small and
statistically indistinguishable from zero ($b=-0.05$, $p=0.46$), a materially different finding from an
earlier version of this analysis that relied on an undercounted Washington dispensary registry and had found
a significant negative effect. We report this null directly rather than omit it. Result 2's per-capita
versus per-dispensary contrast between Colorado, Washington, and Montana remains the paper's evidence for the
license-cap mechanism; this entry-density test was intended as additional, direct evidence on the same
mechanism, and it does not corroborate it once Washington's dispensary data are corrected.

Separately, scaling Result 1's estimated border coefficient ($b=0.341$, a 41 percent premium) by Colorado and
Washington border counties' actual 2025 revenue (\$351 million across 10 counties under the corrected border
definition) implies approximately \$101 million in revenue attributable to the border effect that year. We
present this as an accounting decomposition of an already-estimated coefficient, not as an independent
causal magnitude or a welfare calculation. Higher revenue per dispensary under a binding cap is
observationally consistent with either dispensaries serving a genuinely larger customer base or exploiting
the local market power that restricted entry creates, and \citet{HollenbeckUetake2021} document exactly this
kind of market power for Washington's licensing regime generally; our revenue data contain no prices or
quantities, so we cannot adjudicate between these explanations, and we do not characterize this figure as a
welfare effect.

Converting this revenue figure into an approximate state tax-revenue transfer, Colorado levies a 15 percent
state retail marijuana sales tax on retail transactions (in place of, not in addition to, the standard 2.9
percent general sales tax), while Washington levies a 37 percent excise tax on retail cannabis sales
\citep{CDORMarijuanaSalesTax,WADORCannabisRetailers}. Splitting the \$101 million estimate across the two
states in proportion to each state's own 2025 border-county revenue (\$56.6 million in Colorado, \$44.7
million in Washington) and applying each state's statutory rate implies approximately \$8.5 million in
Colorado state tax revenue and \$16.6 million in Washington state tax revenue -- roughly \$25 million
combined -- attributable to the border premium in 2025. This figure compounds the caveats already noted for
the underlying \$101 million estimate, and additionally assumes the reported sales figures are the pre-tax
base to which each state's statutory rate applies directly; it excludes local-option cannabis taxes and
Colorado's separate wholesale excise tax (levied at the cultivator-to-retailer transfer stage rather than on
retail sales), so it should be read as an illustrative order of magnitude rather than a precise fiscal
accounting.

\newpage
\bibliographystyle{apalike}
\bibliography{references}

\end{document}
